\chapter{Ý tưởng sản phẩm}
\label{ch:y-tuong}

Dùng AI để sinh đề là chuyện đã có nhiều người làm. Chương này nói về những chỗ Kriky quyết khác, và
lý do đằng sau mỗi quyết định.

\section{Luồng tổng thể}
\label{sec:vong-hoc}

Ý tưởng lõi không nằm ở việc sinh câu hỏi, mà ở việc khép kín một vòng tám bước:

\begin{enumerate}
  \item Giáo viên tạo và duyệt đề.
  \item Học sinh làm bài.
  \item Hệ thống đánh giá đáp án và cách tư duy.
  \item Hệ thống xác định lỗi sai hoặc hiểu nhầm khái niệm.
  \item Hệ thống giải thích lỗi, học sinh hỏi lại đến khi hiểu.
  \item Hệ thống tạo câu biến thể của chính câu vừa sai.
  \item Học sinh làm lại, tối đa ba vòng cho mỗi câu.
  \item Mỗi câu chốt ở một trong ba mức điểm, và bài kết thúc.
\end{enumerate}

Hình~\ref{fig:vong-hoc} vẽ tám bước ấy kèm chỗ vòng lặp quay lại, và nó là bản đồ của cả chương
này.

\begin{landscape}
\begin{figure}[p]
  \centering
  \includegraphics[width=\linewidth,height=\textheight,keepaspectratio]{vong-hoc.pdf}
  \caption[Vòng học khép kín]{Vòng học khép kín --- tám bước, và chỗ vòng lặp quay lại}
  \label{fig:vong-hoc}
\end{figure}
\end{landscape}

\section{Hai giai đoạn làm bài của học sinh}
\label{sec:hai-giai-doan}

\emph{Bước 2, 7 và 8 của hình~\ref{fig:vong-hoc}.}

Một bài kiểm tra có hai giai đoạn: giai đoạn 1 là làm bài và nộp, 
giai đoạn 2 là chữa những câu đã sai. Nộp bài kết thúc giai đoạn 1 chứ không kết thúc
bài. Bài chỉ kết thúc khi giai đoạn 2 xong, hoặc khi hạn của giai đoạn 2 hết.

Kéo theo đó là một hệ quả cần nói rõ: một bài kiểm tra có hai trạng thái hoàn thành khác nhau,
\emph{đã nộp} và \emph{đã hoàn thành}, và chúng không thay thế cho nhau được.

\section{Làm sao phân tích lỗi của học sinh?}
\label{sec:nhan-loi}

\emph{Bước 4 của hình~\ref{fig:vong-hoc}.}

Với hình thức trắc nghiệm, việc chấm chỉ là một phép so. Khó khăn nằm ở chẩn đoán lỗi sai của học sinh, 
bởi cùng chọn một phương án sai có thể do nhiều lỗi khác nhau sinh ra. Nếu hệ thống đoán lỗi rồi đi giải 
thích một lỗi học sinh không hề mắc thì nó còn tệ hơn im lặng: em vừa mất thời gian, vừa được dạy sai chỗ.

Vì vậy, mỗi phương án sai phải mang theo lỗi mà nó đại diện, và được soạn cùng câu hỏi
chứ không suy ra lúc chấm. 

Hai luật đi kèm:

\begin{itemize}
  \item Mỗi câu hỏi phải mang theo lời giải có nhiều hơn một cách làm. 
  \item Số phương án không cố định. Một câu có thể có ba, bốn hay năm lựa chọn tuỳ từng câu,
        nhưng chỉ đúng một lựa chọn là đáp án đúng. Mọi phương án còn lại đều là nhiễu, nên đều
        phải có lỗi gắn kèm.
\end{itemize}

\section{Tạo câu hỏi mới dựa trên câu làm sai}
\label{sec:ba-vong}

\emph{Bước 6 và 7 của hình~\ref{fig:vong-hoc}}

Khi học sinh sai câu 3, hệ thống sinh câu 3\`, gọi là câu biến thể. Thứ cần chữa là một lỗi cụ thể vừa xảy ra, không phải một chủ đề.

Câu biến thể giữ nguyên dạng đề và cách làm, chỉ đổi dữ kiện, con số và cách hỏi. Đó là cách duy
nhất kiểm được rằng học sinh đã sửa đúng chỗ mình sai, chứ không phải chỉ nhớ được đáp án.

\section{Thang điểm}
\label{sec:thang-diem}

\emph{Bước 8 của hình~\ref{fig:vong-hoc}.}

\begin{table}[htbp]
\centering
\caption{Ba mức điểm của một câu hỏi}
\label{tab:thang-ba-muc}
\begin{tabular}{@{}cl@{}}
\toprule
\textbf{Điểm} & \textbf{Nghĩa} \\
\midrule
1   & Đúng ngay ở giai đoạn 1 \\
0,5 & Sai ở giai đoạn 1 nhưng chữa được ở giai đoạn 2 \\
0   & Sai và không chữa được \\
\bottomrule
\end{tabular}
\end{table}

Ba mức thay vì hai, vì hai mức sẽ biến việc chữa bài thành công việc không được có ý nghĩa. Một học
sinh sửa được lỗi của mình đã học đúng thứ mà bài kiểm tra muốn dạy, và mức giữa là chỗ duy nhất ghi
nhận điều đó. Vì vậy giai đoạn 2 chỉ cộng chứ không bao giờ trừ: điểm tính ngay sau khi nộp là sàn.

\section{Giáo viên đưa ra quyết định ở đâu?}
\label{sec:ba-cong}

\emph{Bước 1 của hình~\ref{fig:vong-hoc}.}

Với nội dung đề, trợ lí viết còn giáo viên quyết nội dung đó có đến tay học sinh hay không. 
Giáo viên bắt buộc tham gia ở ba điểm, hai ở đầu ra và một ở đầu vào:

\begin{enumerate}
  \item Duyệt đề trước khi phát hành (đầu ra).
  \item Trợ lí hỏi lại khi chưa đủ thông tin, thay vì tự đoán (đầu vào).
\end{enumerate}
